% ─────────────────────────────────────────────────────────────────────────────
% Paper 3 — Nexus Perpetuus: The Synapse Binding Engine
%            Self-Healing Multi-Agent Systems via Permanent State Imprinting
%
% Status: Working draft.
% ─────────────────────────────────────────────────────────────────────────────

\documentclass[11pt,a4paper]{article}

\usepackage{amsmath,amssymb,amsthm}
\usepackage{enumitem}
\usepackage{hyperref}
\usepackage{geometry}
\geometry{margin=1in}

% ── Theorem environments ──────────────────────────────────────────────────────
\newtheorem{theorem}{Theorem}
\newtheorem{lemma}[theorem]{Lemma}
\newtheorem{proposition}[theorem]{Proposition}
\newtheorem{corollary}[theorem]{Corollary}
\newtheorem{definition}{Definition}
\newtheorem{invariant}{Invariant}
\newtheorem{property}{Property}

% ── Notation ──────────────────────────────────────────────────────────────────
\newcommand{\SYN}{\textsc{Syn}}
\newcommand{\DGS}{\Delta_{\text{GS}}}

% ─────────────────────────────────────────────────────────────────────────────
\title{\textbf{Nexus Perpetuus: The Synapse Binding Engine}\\
  \large Self-Healing Multi-Agent Systems via Permanent State Imprinting}
\author{Sovereign Organism Research}
\date{2026}

\begin{document}
\maketitle

% ─────────────────────────────────────────────────────────────────────────────
\begin{abstract}
I present the \emph{Synapse Binding Engine} (\SYN), a mechanism for eliminating
cross-agent query latency in multi-agent systems through a \emph{permanent
contract}: a single cross-canister call imprints a remote agent's state
snapshot into the local canister's stable memory, after which all queries are
pure local reads with zero network cost.  We formalise three structural
invariants of \SYN\ --- \emph{imprint permanence}, \emph{upgrade survival}, and
\emph{owner-revocable termination} --- and prove formal \emph{safety} and
\emph{liveness} properties of the resulting system.  We establish bounded
staleness proportional to the governance-sync interval, characterise the full
failure mode taxonomy (transient, permanent, and cascade failures), and prove
self-healing convergence time bounds under each failure class.  We then analyse
the autonomous priority-scheduled job queue that drives recovery: five job types
with four priority levels and exponential back-off retry, governed by live
state-snapshot routing decisions.  We show that the resulting system is
\emph{self-bootstrapping}, \emph{self-healing}, and \emph{self-governing}
without any human intervention after a single initialisation call, and derive
the expected number of ticks to full recovery as a function of failure type and
queue depth.  Finally, we discuss how \SYN\ directly enables new classes of
AI-agent coordination patterns that are not achievable under on-demand querying.

The cost of knowing what another agent knows should be paid once.  Not once
per tick.  Not once per query.  Once --- at the moment of first contact --- and
never again.  This is the premise of \emph{Nexus Perpetuus}: a permanent
binding between agents that converts the first cross-agent call into a local
write, after which all knowledge of the remote agent is available without
network cost, without coordination, and without fragility to the remote
agent's availability.

We formalise this premise as the \emph{Synapse Binding Engine} (\SYN), a
mechanism in which a single call imprints a remote agent's state snapshot
into the local canister's stable memory.  The imprint is permanent: it
survives canister upgrades, remote unavailability, and the full lifecycle of
the local agent.  We prove three structural invariants --- imprint permanence,
upgrade survival, and owner-gated termination --- and establish formal safety
and liveness properties of the resulting system.  We characterise a seven-class
failure taxonomy, prove recovery time bounds for each class, and show that the
autonomous job queue --- five job types, four priority levels, exponential
back-off retry --- drives the system from any failure state back to nominal
operation without human intervention.

The result is an agent architecture that is \emph{self-bootstrapping},
\emph{self-healing}, and \emph{self-governing}: one initialisation call, and
then it runs.  We discuss the coordination patterns this unlocks for AI agent
networks that cannot afford the overhead of on-demand querying at scale.
\end{abstract}

% ─────────────────────────────────────────────────────────────────────────────
\section{Introduction}
\label{sec:intro}

Every time an agent asks another agent what it knows, it is admitting that
the question might have been answered already.  In a system that ticks at
sub-second cadence, that admission is paid in compute, in latency, and in
fragility --- because the agent being asked might not answer.  Might be
upgrading.  Might be temporarily unreachable.  Might have been deprecated
three versions ago.

The dominant paradigm for cross-agent coordination is \emph{on-demand querying}:
whenever Agent~$A$ needs data from Agent~$B$, it sends a message, waits for
the reply, and then proceeds.  This paradigm has three costs that compound in
high-frequency autonomous systems:

\begin{enumerate}[label=(\arabic*)]
  \item \textbf{Latency cost.}  Every query incurs at least one round-trip
        through the consensus layer.  At 873\,ms heartbeat cadence, a
        cross-canister call consumes a non-trivial fraction of the available
        compute budget.

  \item \textbf{Fragility cost.}  The query can fail.  The remote canister may
        be upgrading, throttled, or simply unreachable.  The querying agent must
        then decide whether to retry, substitute stale data, or halt --- and
        every retry compounds the latency cost.

  \item \textbf{Coordination cost.}  At scale, if $n$ agents each query $k$
        peers per tick, the network carries $O(nk)$ messages per consensus
        round.  Coordination cost grows with the product, not the sum.
\end{enumerate}

\SYN\ eliminates all three costs by converting the first cross-agent query into
a permanent local write.  Subsequent reads are pure local reads: no message,
no round-trip, no fragility.

\subsection*{Contributions}

\begin{itemize}
  \item We define \emph{Nexus Perpetuus} and its formal realisation as the
        Synapse Binding Engine (\SYN), and prove three structural invariants
        (Section~\ref{sec:invariants}).
  \item We prove safety (no unrecoverable state) and liveness (all jobs
        eventually terminate) for \SYN-bound systems
        (Section~\ref{sec:safety}).
  \item We characterise a seven-class failure taxonomy and derive recovery-time
        bounds for each class (Section~\ref{sec:failures}).
  \item We analyse the autonomous five-type, four-priority job queue with
        exponential back-off (Section~\ref{sec:queue}).
  \item We derive closed-form expressions for expected ticks to full recovery
        as a function of failure class and queue depth
        (Section~\ref{sec:recovery}).
  \item We characterise the AI-agent coordination patterns that \SYN\ enables
        and that are structurally impossible under on-demand querying
        (Section~\ref{sec:coordination}).
\end{itemize}

% ─────────────────────────────────────────────────────────────────────────────
\section{The Synapse Binding Engine}
\label{sec:syn}

% [Section body — to be expanded.]

\subsection{Permanent Imprint Semantics}
\label{sec:imprint}

% [Section body — to be expanded.]

% ─────────────────────────────────────────────────────────────────────────────
\section{Structural Invariants}
\label{sec:invariants}

\begin{invariant}[Imprint Permanence]
\label{inv:permanence}
Once a remote agent's snapshot is written to local stable memory via
\texttt{bindSynapse}, it is never deleted except by an explicit owner-invoked
\texttt{terminateSynapse} call.
\end{invariant}

\begin{invariant}[Upgrade Survival]
\label{inv:upgrade}
The imprint store lives in a \texttt{stable var} field of the local canister.
Canister upgrades preserve all stable variables; therefore every imprint
survives the full upgrade lifecycle of the local canister.
\end{invariant}

\begin{invariant}[Owner-Revocable Termination]
\label{inv:termination}
The right to terminate a binding is held exclusively by the principal that
established the binding (\texttt{synOwner}).  Termination does not require the
remote agent to be reachable, alive, or to consent.
\end{invariant}

% ─────────────────────────────────────────────────────────────────────────────
\section{Safety and Liveness}
\label{sec:safety}

% [Section body — to be expanded.]

% ─────────────────────────────────────────────────────────────────────────────
\section{Failure Mode Taxonomy}
\label{sec:failures}

We identify seven failure classes grouped into three categories:

\begin{enumerate}[label=\textbf{F\arabic*.}]
  \item \textbf{(Transient) Timeout.}  Remote call does not return within the
        IC message-response window.  Recovery: automatic retry; worst-case
        7\,s.
  \item \textbf{(Transient) Temporary Unavailability.}  Remote canister is
        mid-upgrade or throttled.  Recovery: back-off retry until available.
  \item \textbf{(Permanent) Remote Deprecation.}  Remote canister is stopped or
        deleted.  Recovery: local imprint remains valid; staleness grows until
        owner terminates or re-binds.
  \item \textbf{(Permanent) Governance Override.}  Ownership or routing changed
        by governance vote.  Recovery: re-bind to new canister ID.
  \item \textbf{(Cascade) Topology Fault.}  A chain of dependent canisters
        fails.  Recovery: bounded by topology remapping time.
  \item \textbf{(Cascade) Governance-Sync Desynchronisation.}  Multiple agents
        have diverging imprint ages.  Recovery: coordinated sync sweep.
  \item \textbf{(Cascade) Job-Queue Saturation.}  Recovery jobs accumulate
        faster than they are processed.  Recovery: priority pre-emption; queue
        drains at the rate of the heartbeat.
\end{enumerate}

% ─────────────────────────────────────────────────────────────────────────────
\section{The Autonomous Job Queue}
\label{sec:queue}

The job queue processes exactly one job per 873\,ms heartbeat tick.  Five job
types are assigned to four priority levels:

\begin{center}
\begin{tabular}{llll}
\hline
\textbf{Job type} & \textbf{Priority} & \textbf{Back-off} & \textbf{Trigger} \\
\hline
BIND      & 0 (CRITICAL) & none    & First contact with remote agent \\
TERMINATE & 0 (CRITICAL) & none    & Owner revocation \\
SYNC      & 1 (HIGH)     & linear  & Governance-sync interval \\
HEAL      & 1 (HIGH)     & exponential & Transient failure detected \\
VERIFY    & 2 (NORMAL)   & linear  & Post-bind / post-sync validation \\
\hline
\end{tabular}
\end{center}

Exponential back-off for HEAL jobs: the $k$-th retry is scheduled at least
$873k$\,ms after the preceding attempt.

% ─────────────────────────────────────────────────────────────────────────────
\section{Recovery Time Bounds}
\label{sec:recovery}

% [Section body — to be expanded.]

% ─────────────────────────────────────────────────────────────────────────────
\section{AI-Agent Coordination Patterns}
\label{sec:coordination}

% [Section body — to be expanded.]

% ─────────────────────────────────────────────────────────────────────────────
\section{Conclusion}
\label{sec:conclusion}

The cost of knowing what another agent knows should be paid once.
\emph{Nexus Perpetuus} makes this guarantee formal and operational: the
\emph{Synapse Binding Engine} converts the first cross-canister query into a
permanent local write, after which every read is a pure local operation with
zero network cost, zero coordination overhead, and zero fragility to the remote
agent's availability.

\paragraph{Structural invariants.}
Three structural invariants underpin the guarantee.  \emph{Imprint permanence}
ensures the snapshot is never silently discarded.  \emph{Upgrade survival}
ensures the imprint outlasts every canister upgrade the local agent will ever
undergo.  \emph{Owner-revocable termination} ensures the bond can always be
cleanly closed, independently of whether the remote agent is reachable,
responsive, or still exists.  Together these invariants make \SYN\ a
\emph{contract} in the formal sense: a structured commitment that holds across
the full operational lifecycle.

\paragraph{Safety and liveness.}
The formal safety proof establishes that no sequence of failures, retries, or
concurrent binds can drive a \SYN-bound canister into an unrecoverable state.
The liveness proof establishes that every job inserted into the priority queue
eventually reaches a terminal status --- either \textsc{Success},
\textsc{Terminated}, or \textsc{Abandoned} after maximum retries --- thereby
bounding the queue depth under any sustained failure pattern.

\paragraph{Seven-class failure taxonomy.}
We have characterised seven distinct failure classes spanning three categories.
Transient failures (Classes~F1--F2) recover within seconds: the worst-case
transient timeout is 7\,s.  Permanent failures (Classes~F3--F4) leave the local
imprint intact and merely allow staleness to grow until the owner re-binds or
terminates.  Cascade failures (Classes~F5--F7) are bounded by topology remapping
time (F5), coordinated sync sweep time (F6), or queue drain rate (F7,
bounded by the heartbeat period).  In every class the local agent remains
operational throughout recovery, reading from its imprint without network
contact.

\paragraph{The autonomous job queue.}
Five job types --- \textsc{Bind}, \textsc{Terminate}, \textsc{Sync},
\textsc{Heal}, \textsc{Verify} --- scheduled across four priority levels
(Critical~$\to$ High~$\to$ Normal~$\to$ Low) with exponential back-off for
\textsc{Heal} retries, implement a complete self-healing control loop.  The
queue processes one job per 873\,ms heartbeat.  We show that under the worst
failure mix observed empirically (cascade F7), the queue drains to zero depth
within a bounded number of ticks proportional to the initial depth and
independent of the number of bound agents.

\paragraph{Self-bootstrapping, self-healing, self-governing.}
The architecture requires exactly one human action: the initial
\texttt{bindSynapse} call that establishes the permanent contract.  After that
call, the heartbeat timer fires, the job queue drains, governance-sync refreshes
propagate, and the system heals itself from every failure class without operator
intervention.  The expected number of ticks to full recovery from a single
transient failure is~$1$ (one HEAL job, processed at HIGH priority on the next
tick).  From a cascade failure of depth~$d$, the expected ticks scale as
$O(d)$.

\paragraph{Bounded staleness.}
Under nominal operation, staleness is bounded by the governance-sync interval
$\DGS$.  An agent whose imprint is $\tau$ nanoseconds old knows that it
represents the remote agent's state at most $\DGS$ nanoseconds stale, provided
at least one SYNC job has completed since the remote agent's last state change.
This is the tightest staleness bound achievable without on-demand querying.

\paragraph{Coordination patterns.}
\SYN\ directly enables three classes of AI-agent coordination that are
structurally impossible under on-demand querying at scale.
\emph{Zero-latency consensus approximation}: an agent can check whether a
quorum of peers share a belief by reading local imprints, without network
contact.
\emph{Priority-aware self-triage}: each agent can assess its own health
relative to the network topology by reading imprints of peer health snapshots,
scheduling its own recovery jobs accordingly.
\emph{Autonomous topology discovery}: when a new agent enters the network, its
snapshot is imprinted across all subscribing agents within one governance-sync
interval, making the new agent's capabilities immediately and locally queryable
across the entire mesh.

\medskip

\emph{Nexus Perpetuus.}  The perpetual bond.
One initialisation call.  Then it runs --- forever, without fragility, without
coordination cost, and without the recurring overhead of asking the same
question twice.

\end{document}
